% =====================================================================
% 第二部分 · 中国近现代史纲要
% =====================================================================

\zzchapter{第二部分\quad 中国近现代史纲要}

\zzsection{一、旧民主主义革命阶段}

\zzsub{清末新政}

\zzcard{清末新政}{%
  \kw{“撤绿立新”}：\bk{裁撤绿营，建立新军}。
}

\zzsub{同盟会起义}

\zzcard{同盟会起义}{%
  \kw{“大黄花”}：同盟会发动的\bk{影响最大}的起义——\kw{黄花岗起义}（又称广州起义、广州三二九起义）。\\
  \zzwarn{同盟会成立后发动的\bk{第一次}武装起义是\kw{萍浏醴起义}。}\par\addvspace{2pt}
  \zzwarn{三元里人民的抗英斗争是中国近代史上中国人民\bk{第一次大规模的反侵略武装斗争}。}
}

\zzsub{戊戌维新}

\zzcard{戊戌维新}{%
  戊戌维新运动是\kw{中国民族资产阶级登上政治舞台的第一次表演}。
}

\zzsub{林则徐 严复 魏源 郑观应}

\zzcard{近代开眼看世界的人物}{%
  \kw{“徐四睁眼见阎王喂海狮，郑观应是公主”}\\
  \begin{zztable2}
    人物 & 著作与主张 \\
    \midrule
    林则徐 & 《四洲志》，\bk{睁眼看世界第一人} \\
    严复 & \bk{救亡图存} \\
    魏源 & 《海国图志》，\bk{“师夷长技以制夷”} \\
    郑观应 & 《盛世危言》，\bk{“君民共主”} \\
  \end{zztable2}
}

\zzsub{思想交锋}

\zzcard{思想交锋}{%
  \bk{维新派与守旧派的论战}是\kw{近代中国资产阶级思想与封建主义思想的第一次正面交锋}。\\
  论战所涉及的领域十分广泛，进一步开阔了新型知识分子的眼界，解放了人们长期受到束缚的思想。通过论战，西方资产阶级社会政治学说在中国得到进一步传播，戊戌变法运动的帷幕随之拉开。\\
  \bk{维新派与洋务派分歧的实质}是\kw{要不要学习西方的政治制度与思想文化}。\\
  \zzwarn{维新派\bk{技术思想都学}，洋务派\bk{只学技术不学思想}。}
}

\zzsection{二、新文化运动及五四运动}

\zzsub{中国早期信仰马克思主义的人物的队伍}

\zzcard{中国早期信仰马克思主义的人物的队伍}{%
  \kw{“骨干是主体”}：队伍主要由三种人组成，其\bk{主体}部分是\kw{五四爱国运动的左翼骨干}。
}

\zzsub{新文化运动}

\zzcard{新文化运动}{%
  五四以前的\bk{新文化运动}的倡导者就已经把\kw{改造国民性}置于优先地位。新文化运动的倡导者：\\
  ①以\bk{进化论观点和个性解放思想}为主要武器；\\
  ②\bk{猛烈抨击以孔子为代表的“往圣先贤”}（通过批判孔学，动摇了封建正统思想的统治地位，在中国社会掀起一股思想解放的潮流）；\\
  ③批判孔学，是为了\bk{给中国发展资本主义扫清障碍}（屁孔发紫）；\\
  ④大力提倡\bk{新道德、反对旧道德}，提倡\bk{新文学、反对旧文学}，包括提倡白话文、反对文言文；提倡\bk{民主、反对专制}，提倡\bk{科学、反对迷信盲从}。\\
  \zzwarn{新文化运动时期，中国\bk{还没有马克思主义}，后来十月革命给中国送来了马克思主义。五四以前新文化运动的领导人物，还没有马克思主义的批判精神，因此看问题很片面。}
}

\zzsub{李大钊}

\zzcard{李大钊}{%
  ①\kw{率先在中国大地上举起马克思主义旗帜}；\\
  ②《法俄革命之比较观》向中国人民\bk{第一次正确地阐述了十月革命的性质}；\\
  ③《庶民的胜利》《布尔什维主义的胜利》深刻揭露了第一次世界大战的本质，\bk{热情歌颂}了十月革命和布尔什维主义的胜利；\\
  ④《我的马克思主义观》明确地把马克思主义称为“世界改造原动的学说”，并对马克思的唯物史观、剩余价值学说和阶级斗争理论作了比较系统的介绍，这表明\kw{李大钊已经成为中国的第一个马克思主义者}。\\
  《我的马克思主义观》也\bk{标志着马克思主义在中国进入一个比较系统的传播阶段}。
}

\zzsection{三、帽子题及易混对比}

\zzsub{土地政策}

\zzcard{土地政策}{%
  \kw{“康熙底线心存贞洁”}\\
  \begin{zztable2}
    时期 & 土地政策 \\
    \midrule
    土地革命时期 & 限制富农 \\
    抗日战争时期 & 减租减息 \\
    解放战争时期 & 征收富农多余土地 \\
    新中国成立 & 保存富农经济 \\
  \end{zztable2}
}

\zzsub{精神辨析}

\zzcard{精神辨析}{%
  \bk{五四精神}：爱国、进步、民主、科学（五四精神的\bk{核心是爱国主义}）。\\
  \bk{黄埔精神}：爱国、革命。
}

\zzsub{工人罢工运动}

\zzcard{工人罢工运动}{%
  \kw{“慕容云海是起点，张翰是终点；六胜五不胜，省港最长”}\\
  \begin{zztable2}
    罢工 & 地位 \\
    \midrule
    香港海员罢工 & \bk{起点} \\
    京汉铁路大罢工 & \bk{终点（高潮）} \\
    安源路矿工人罢工 & \bk{完全胜利} \\
    开滦五矿工人罢工 & \bk{非完全胜利} \\
    省港大罢工 & \bk{持续时间最长} \\
  \end{zztable2}
}

\zzsub{中心问题与台儿庄大捷}

\zzcard{中心问题}{%
  \bk{中国革命的中心问题}是\kw{无产阶级领导权}。\\
  \bk{国民革命的中心问题}是\kw{农民问题}。
}

\zzcard{台儿庄大捷}{%
  \kw{“正态分布”}：抗战爆发后中国\bk{正面战场}取得的\kw{首次重大胜利}——台儿庄大捷（国民党正面战场，1938 年）。
}

\zzsub{秋收起义与南昌起义}

\zzcard{秋收起义与南昌起义}{%
  \zzsub{区别}
  ★\bk{秋收起义}：【\kw{工农革命军}】的旗帜；\\
  ★\bk{南昌起义}：【\kw{国民革命军}】的旗帜；\\
  ★\bk{秋收起义不仅是军队的行动}，而且有数量众多的\bk{工农武装}参加。\\
  \zzsub{联系}
  ★\bk{秋收起义的目标是攻打中心城市长沙}。在攻打中心城市受挫后，毛泽东果断改变计划，决定向敌人控制比较薄弱的农村转移——这一点与南昌起义是相同的。\\
  ★南昌起义和秋收起义都是\bk{中国共产党独立领导的革命战争}。
}

\zzsub{延安整风运动}

\zzcard{延安整风运动}{%
  \kw{“尊师重教”}：延安整风运动的\bk{重点是教条主义}。\\
  \zzwarn{整风运动\bk{最主要的任务}：反对主观主义。整风运动中\bk{第一个重要的问题}：学风问题。}
}

\zzsub{统一战线辨析}

\zzcard{统一战线辨析}{%
  \bk{工农民主统一战线}：大革命后因汪蒋叛变，民族小资产阶级、国民党都走了，只剩下工人和农民留在共产党，故此时是工农民主统一战线。\\
  \bk{人民民主统一战线}：解放战争时期我们争取尽可能多的人民（除工农外也包括富农、民族资产阶级等）与国民党进行最终决战，此时形成人民民主统一战线。
}

\zzsub{维新派和辛亥革命对比}

\zzcard{维新派和辛亥革命对比}{%
  \begin{zztable2}
    维新派 & 辛亥革命 \\
    \midrule
    有利于\bk{民主思想}在中国的传播，有利于人们的思想解放 & 结束了中国延续两千多年的\bk{封建君主专制制度} \\
    把顽固的\bk{封建主义思想壁垒}打开了一个缺口 & 建立了中国历史上第一个\bk{资产阶级共和政府} \\
    形成了广泛的\bk{文化革新运动} & 使\bk{民主共和的观念}开始深入人心 \\
    以维新运动为\bk{起点}，资产阶级新文化开始打破封建文化独占文化阵地的局面 & \\
    对中国近代教育的发展起了积极的推动作用 & \\
  \end{zztable2}
}

\zzsub{党对军队的绝对领导}

\zzcard{党对军队的绝对领导}{%
  \kw{“Fa男鸡窝股鼓”}\\
  \bk{发端于南昌起义}、\bk{奠基于三湾改编}、\bk{定型于古田会议}。\\[2pt]
  ①\bk{古田会议}：确立了马克思主义建党建军原则，\kw{【思想建党，政治建军】}，创造性地解决了在农村环境中、在党组织和军队以农民为主要成分的条件下，如何保持党的无产阶级先锋队性质和建设党领导的新型人民军队的重大问题。（古田先锋军队）\\
  ②\bk{三湾改编}：确立了\kw{支部建在连上}的原则和制度。建设无产阶级领导的\bk{新}型人民军队的\bk{重要开端}：三湾改编。（三县）
}

\zzsub{新旧三民主义}

\zzcard{新旧三民主义}{%
  \begin{zztable2}
    旧三民主义 & 新三民主义 \\
    \midrule
    驱除鞑虏，恢复中华 & \bk{民族}：突出\bk{反帝}内容，对外实行民族独立，对内主张民族平等 \\
    创立民国 & \bk{民权}：民主权利“为一般平民所共有”，不应为“少数人所得私” \\
    平均地权 & \bk{民生}：平均地权和\bk{节约资本}（后又提出耕者有其田），并提出改善工农生活状况 \\
  \end{zztable2}
  \zzwarn{“三民无主”：新旧三民主义都是\bk{民族、民生、民权}，\bk{从来没有“民主”}，有民主就错。}\par\addvspace{2pt}
  \zzkey{新三民主义的政纲同中国共产党在民主革命阶段的纲领基本一致，因而有了\bk{国共合作的政治基础}。}
}

\zzsection{四、旧民主主义革命的终结和新民主主义革命的开端}

\zzsub{革命性质与领导权}

\zzcard{旧民主主义革命的终结和新民主主义革命的开端}{%
  ①\bk{护法运动的失败}不仅是孙中山个人的失败，也标志着整个中国民族资产阶级领导的\kw{旧民主主义革命的终结}；\\
  ②\bk{五四运动}成为中国革命的新阶段即\kw{新民主主义革命阶段的开端}。\\[2pt]
  \zzwarn{无产阶级的领导权是中国革命的中心问题，也是新民主主义革命理论的核心问题。\bk{区别新、旧两种不同范畴的民主主义革命的根本标志是，革命的领导权是掌握在无产阶级手中还是掌握在资产阶级手中。}}\par\addvspace{2pt}
  \zzwarn{注意：虽然“新民主主义革命”有个“新”字，但其\bk{性质依然是资产阶级民主主义革命}（由社会性质、革命对象、革命任务等决定）。但是，中国革命已不是旧式的、一般的资产阶级民主主义革命，而是\bk{新式的、特殊的}资产阶级民主主义革命。}
}

\zzsub{转折点}

\zzlongcard{转折点}{%
  \kw{“1947 米脂县，革命两个转折点”}\\
  ①\bk{1947 年 12 月}，中共中央在陕北\bk{米脂县杨家沟}召开会议，制定了夺取全国胜利的行动纲领。毛泽东指出，中国革命已经发展到了一个历史的转折点：\\
  \hspace{1em}→\bk{蒋介石 20 年反革命统治由发展到消灭的转折点}；\\
  \hspace{1em}→\bk{一百多年来帝国主义在中国的统治由发展到消灭的转折点}。\\
  ②\bk{1949 年 3 月 25 日}，毛泽东等中央领导人与中央机关、人民解放军总部\bk{进驻北平}香山，标志着\kw{中国革命重心从农村转向城市}。\\[2pt]
  \zzwarn{辨析：\bk{秋收起义}——城市转农村；\bk{进驻香山}——农村转城市。}\par\addvspace{2pt}
  \zzwarn{《中共中央关于党的百年奋斗重大成就和历史经验的决议》指出，从进攻大城市转为向农村进军，是\bk{中国革命具有决定意义的新起点}。}\par\addvspace{2pt}
  \zzwarn{\bk{长征}实现了中国共产党和中国革命事业从挫折走向胜利的伟大转折。}\par\addvspace{2pt}
  \zzwarn{\bk{抗日战争的伟大胜利}，是近代以来中国人民反抗外敌入侵第一次取得完全胜利的民族解放斗争，是中华民族从近代以来陷入深重危机走向伟大复兴的历史转折点。}
}

\zzsection{五、毛泽东思想形成与发展各时期著作}

\zzsub{四阶段总表}

\zzcard{毛泽东思想形成与发展各时期著作}{%
  \begin{zztable2}
    阶段 & 著作 \\
    \midrule
    \bk{①开始萌芽}：一阶级两农民，新民主主义提出 & 《中国社会各阶级分析》《国民革命与农民运动》《湖南农民运动考察报告》 \\
    \bk{②开始形成}：红星冈本，农村包围城市、武装夺取政权（土地革命时期） & 《中国的红色政权为什么能够存在》《星星之火，可以燎原》《井冈山的斗争》《反对本本主义》 \\
    \bk{③趋于成熟}：三论两党改造政府，新民主主义系统阐述 & 《实践论》《矛盾论》《〈共产党人〉发刊词》《中国革命和中国共产党》《新民主主义论》《改造我们的学习》《论联合政府》 \\
    \bk{④继续发展}：两论两在一关于 & 《在中国共产党第七届中央委员会第二次全体会议上的报告》《在晋绥干部会议上的讲话》《论人民民主专政》《论十大关系》《关于正确处理人民内部矛盾的问题》 \\
  \end{zztable2}
  \zzkey{毛泽东思想命名：\kw{中共七大}。}
}

\zzsub{重要著作要点}

\zzcard{《湖南农民运动考察报告》}{%
  总结\bk{辛亥革命失败原因}：\kw{国民革命需要一个大的农村变动}。
}

\zzcard{《中国社会各阶级分析》}{%
  ①分析了各阶级在革命中的作用和地位；\\
  ②提出\kw{“谁是我们的敌人？谁是我们的朋友”}这个革命的\bk{首要问题}。
}

\zzcard{《反对本本主义》}{%
  ①第一次提出\kw{没有调查就没有发言权}；\\
  ②第一次提出\bk{中国革命斗争的胜利靠中国同志了解中国情况}，标志\bk{党独立自主的思想初步形成}；\\
  ③基本形成了\bk{毛泽东思想}。
}

\zzcard{《论联合政府》}{%
  1945 年，毛泽东在党的七大所作的政治报告《论联合政府》中，进一步把新民主主义的政治、经济和文化与党的基本纲领联系起来，进行了具体阐述，形成了\bk{新民主主义的基本纲领}。\\
  \zzsub{新民主主义的经济纲领}
  ①没收封建地主阶级的土地归农民；\\
  ②没收官僚资本归新民主主义国家；\\
  ③保护民族工商业（也是新民主主义经济纲领中极具特色的一项内容）。
}

\zzcard{《论十大关系》}{%
  ①\bk{一个基本方针}：\kw{调动一切积极因素为社会主义事业服务}；\\
  ②明确提出\kw{以苏为鉴}，独立自主地探索适合中国情况的社会主义建设道路（另一说法：明确了建设社会主义的根本思想是\bk{根据本国情况走自己的路}）；\\
  ③\bk{标志着党探索中国社会主义建设道路的良好开端}（十全十美好开端）；\\
  ④\bk{第一大关系也是最重要的关系}：\kw{重工业、轻工业和农业的关系}。
}

\zzlongcard{《关于正确处理人民内部矛盾的问题》}{%
  \zzsub{具体方针}
  ①\bk{政治思想领域}：实行“\kw{团结—批评—团结}”；\\
  ②\bk{党与民主党派的关系}：实行“\kw{长期共存、互相监督}”；\\
  ③\bk{科学文化工作}：实行“\kw{百花齐放、百家争鸣}”；\\
  ④\bk{物质利益分配}及其他各项工作：实行“\kw{统筹兼顾、适当安排}”；\\
  ⑤\bk{人民群众和政府机关}：坚持“\kw{民主集中制原则}”；\\
  ⑥\bk{民族之间}：实行“\kw{民主平等，团结互助}”。\\
  \zzwarn{辨析：“惩前毖后、治病救人”的方针，是中国共产党在\bk{延安整风运动}中提出的对待犯错误同志的正确方针。}\par\addvspace{2pt}
  在该报告中，毛泽东明确提出了\bk{中国工业化道路}的问题，主要是指\bk{重工业和轻工业、农业的发展关系问题}，强调要走一条有别于苏联的中国工业化道路。\\
  走中国工业化道路的思想，是党探索我国社会主义建设道路的一个重要思想，强调要正确处理重工业和轻工业、农业的关系，符合中国人口多、工业基础薄弱的实际。\\
  \bk{内容}：以农业为基础、工业为主导，\kw{农轻重为序}发展国民经济的总方针，以及一整套\kw{“两条腿走路”}的工业化发展思路。
}

\zzcard{《论人民民主专政》}{%
  《论人民民主专政》是毛泽东在中国新民主主义革命取得决定性胜利、全国性政权即将建立的时刻，为纪念中国共产党成立 28 周年而撰写的文章。\\
  文章论述了即将成立的中华人民共和国的国家性质，各阶级在国家中的地位及其相互关系，国家对内、对外政策等，\bk{奠定了新中国国家政权的理论基础和一定发展阶段上的政策基础}，为即将成立的新中国作了\bk{政治理论准备}。\\
  \zzkey{中共七届二中全会的决议和毛泽东的《论人民民主专政》，构成了《中国人民政治协商会议共同纲领》的基础。}
}

\zzsection{六、三大辨析}

\zzsub{党的三大优良作风}

\zzcard{党的三大优良作风}{%
  \kw{“实群批”}：\bk{理论联系实际}、\bk{密切联系群众}、\bk{批评与自我批评}。\\
  \zzwarn{1945 年七大提出三大优良作风。}
}

\zzsub{新民主主义革命三大法宝}

\zzcard{新民主主义革命三大法宝}{%
  \kw{“战斗党”}：\bk{统一战线}、\bk{武装斗争}、\bk{党的建设}。\\
  \zzwarn{1939 年《〈共产党人〉发刊词》提出三大法宝。}
}

\zzsub{统一战线}

\zzcard{统一战线的必要性与可能性}{%
  中国革命\bk{必须}建立最广泛的统一战线，这是由：\\
  ①\bk{中国半殖民地半封建社会的阶级状况}所决定的；\\
  ②\bk{中国革命的长期性、残酷性及其发展的不平衡性}所决定的。\\
  \zzwarn{辨析：半殖民地半封建的中国社会交织在一起的诸多矛盾，决定的是建立最广泛的统一战线的\bk{可能性}，而不是\bk{必要性}。决定\bk{必要性}的是\bk{阶级、不平衡}。\kw{必须选国情，可能选矛盾}。}
}

\zzcard{抗日民族统一战线与“三三制”原则}{%
  ①\bk{策略总方针}：\kw{发展进步势力，争取中间势力，孤立顽固势力}。\\
  \hspace{1em}\bk{进步势力}：工人、农民和城市小资产阶级；\\
  \hspace{1em}\bk{中间势力}：民族资产阶级、开明绅士和地方实力派；\\
  \hspace{1em}\bk{顽固势力}：大地主、大资产阶级的抗日派，即以蒋介石集团为代表的国民党亲英美派。\\
  ②\bk{“三三制”原则}：共产党员、党外进步人士和中间派\bk{各占 1/3}。\\
  \hspace{1em}共产党员：代表工人和农民；党外进步人士：代表小资产阶级；\\
  \hspace{1em}\bk{中间派}：代表民族资产阶级和开明绅士（\bk{无地方实力派}）。
}

\zzsub{工农武装割据}

\zzcard{工农武装割据}{%
  \kw{“武弟兄”}：\bk{土地革命}、\bk{武装斗争}、\bk{根据地建设}。\\
  中国革命走农村包围城市、武装夺取政权的道路，根本在于处理好三者之间的关系：\\
  \bk{土地革命}是中国革命的\kw{基本内容}；\\
  \bk{武装斗争}是中国革命的\kw{主要形式}，是农村根据地建设和土地革命的强有力保证；\\
  \bk{农村革命根据地}是中国革命的\kw{战略阵地}，是进行武装斗争和开展土地革命的依托。\\
  在党的领导下，实现了三者的密切结合和有机统一。\\
  \zzwarn{“发展进步势力、争取中间势力、孤立顽固势力”是\bk{“巩固抗日民族统一战线的策略总方针”}。}
}

\zzsub{毛泽东思想活的灵魂}

\zzlongcard{毛泽东思想活的灵魂}{%
  \kw{“是群主”}：\bk{实事求是}、\bk{群众路线}、\bk{独立自主}。\\
  \zzsub{实事求是}
  是\bk{毛泽东思想的根本点}、\bk{毛泽东思想的精髓}、\bk{党的思想路线的核心}、\bk{党的根本思想路线}。\\
  \zzsub{群众路线}
  是\bk{我们党的生命线和根本工作路线}，是\bk{我们党永葆青春活力和战斗力的重要传家宝}。\\
  \zzsub{独立自主}
  是\bk{中华民族的优秀传统}、\bk{中华民族精神之魂}、\bk{中国共产党和中华人民共和国立党立国的重要原则}。\\
  \bk{独立自主的探索和实践精神}是\bk{我们党全部理论和实践的立足点}，是\bk{党和人民事业不断从胜利走向胜利的根本保证}。\\
  \zzsub{走自己的路}
  是\bk{党的全部理论和实践立足点}，是\bk{党百年奋斗得出的历史结论}。\\
  \zzsub{中华优秀传统文化}
  是\bk{中华民族的根和魂}、\bk{最深厚的文化软实力}、\bk{中国特色社会主义植根的沃土}、\bk{我们在世界文化激荡中站稳脚跟的根基}。
}

\zzsection{七、重要专题}

\zzsub{马克思主义中国化}

\zzlongcard{马克思主义中国化}{%
  \zzsub{三次飞跃}
  ①\bk{中共七大}确立\kw{毛泽东思想}为党的指导思想（\bk{第一次}）；\\
  ②\bk{中国特色社会主义理论体系}，改革开放后（\bk{第二次}）；\\
  ③\bk{习近平新时代中国特色社会主义思想}（\kw{新的飞跃，非第三次}）。\\
  \zzsub{两次结合}
  ①马克思主义同\bk{中国革命实际}的结合——正式提出命题于\kw{六届六中全会《论新阶段》}（\zzwarn{共产党成立就开始了}）；\\
  ②马克思主义同\bk{中国社会主义建设实际}的结合——\kw{1956 年《论十大关系》}。\\
  \zzsub{两个结合}
  ①把马克思主义基本原理同\bk{中国具体实际}相结合；\\
  ②把马克思主义基本原理同\bk{中华优秀传统文化}相结合。\\[2pt]
  \kw{“六六结合、七大飞跃、五六结合、改革飞跃”}\\
  \zzwarn{毛泽东思想\bk{无改革}；中国特色社会主义理论体系\bk{无革命}；马克思主义中国化时代化\bk{都有}。}\par\addvspace{2pt}
  \zzwarn{魂脉——\bk{马克思主义}；根脉——\bk{中华优秀传统文化}。}
}

\zzsub{我国的主要矛盾改变}

\zzcard{我国的主要矛盾改变}{%
  \begin{zztable2}
    阶段 & 主要矛盾 \\
    \midrule
    \bk{社会主义建设时期}（毛）八大 & 人民对于经济文化迅速发展的需要同当前经济文化不能满足人民需要的状况之间的矛盾（\kw{经济}）；人民对建立先进工业国的要求同落后的农业国现实之间的矛盾 \\
    \bk{改革开放后}（邓）十一届六中《决议》 & 人民日益增长的物质文化需要同落后的社会生产之间的矛盾（\kw{物质}） \\
    \bk{新时代}（习）十九大 & 人民日益增长的美好生活需要和不平衡不充分的发展之间的矛盾（\kw{美好生活}） \\
  \end{zztable2}
}

\zzsub{关于思想路线}

\zzlongcard{关于思想路线}{%
  ①\bk{《实践论》《矛盾论》}：从马克思主义认识论的新高度总结中国共产党的历史经验，揭露和批判党内的主观主义，尤其是教条主义错误，深入论证马克思列宁主义基本原理同中国具体实际相结合的原则，\kw{科学阐明了党的马克思主义的思想路线}。\\
  ②\bk{延安整风}：一切从实际出发、理论联系实际、\kw{实事求是}的马克思主义\bk{思想路线}在全党范围\bk{确立}。\\
  ③\bk{关于真理标准问题的讨论}：是继五四运动和延安整风运动之后又一场马克思主义思想解放运动，其实质在于\bk{是不是坚持马克思列宁主义、毛泽东思想}。讨论\kw{明确地解决了党的思想路线问题}，成为\bk{拨乱反正和改革开放的思想先导}，为党重新确立实事求是思想路线、纠正长期以来的“左”倾错误、实现历史的转折\bk{作了思想理论准备}。\\
  ④\bk{十一届三中全会}：冲破长期“左”的错误的严重束缚，彻底否定了“两个凡是”的错误方针，高度评价了关于真理标准问题的讨论，断然否定“以阶级斗争为纲”的指导思想，作出把工作重点转移到社会主义现代化建设上来和实行改革开放的战略决策，\kw{重新确立了马克思主义的思想路线、政治路线和组织路线}。\\
  \hspace{1em}另一说法：重新确立\bk{解放思想、实事求是}思想路线。\\
  \hspace{1em}全会恢复了党的民主集中制的优良传统，审查解决了历史上遗留的一批重大问题和一些重要领导人的功过是非问题；形成了\bk{以邓小平为核心的党的中央领导集体}，\bk{揭开了社会主义改革开放的序幕}。以这次全会为起点，\kw{中国进入了改革开放和社会主义现代化建设的历史新时期}。
}

\zzsub{会议：1921—1978}

\zzcard{会议顺口溜}{%
  \zzkey{一大立党二大纲，三大国共要合作，国民一大合作成，四大领导工农盟，五大纠正陈独秀，八七枪杆出政权，古田军队要增强，遵义军事与组织，瓦窑堡一起来抗日，洛川全面来抗战，六届六中马中化，六届七中整完风，七大确定毛领导。}
}

\zzlongcard{会议要点详解}{%
  \zzsub{八七会议}
  总结了大革命\bk{失败}的教训（巴适），确定了\bk{土地革命和武装起义}的方针。\\
  \zzsub{遵义会议（遵君嘱）}
  ①遵义会议的一系列重大决策，是在中国共产党同共产国际的联系\bk{中断}的情况下\kw{独立自主}地作出的；\\
  ②遵义会议开启了党\kw{独立自主}解决中国革命实际问题的新阶段；\\
  ③遵义会议的\bk{鲜明特点}是\kw{坚持真理、修正错误}（遵义修理）。\\
  \zzsub{瓦窑堡会议}
  解决党的\bk{政治路线}问题（瓦罗兰，统一战线）。\\
  \zzsub{洛川会议}
  制定了\bk{抗日救国十大纲领}，强调要打倒日本帝国主义，关键在于：\kw{使已然发动的抗战成为全面的全民族的抗战}。\\
  \zzsub{六届六中全会}
  ①毛泽东明确提出\kw{“马克思主义的中国化”}命题；②全会基本纠正了\bk{王明的右倾错误}。\\
  \zzsub{六届七中全会}
  通过了《关于若干历史问题的决议》（1945 年），对党的若干重大历史问题作出结论，使全党特别是党的高级干部对中国革命基本问题的认识达到了一致，增强了全党团结，充分肯定了确立毛泽东在全党领导地位的重大意义，为七大的胜利召开在思想上做了充分准备，\kw{标志着延安整风运动的胜利结束}。\\
  \zzsub{中共七大}
  ①把\kw{毛泽东思想}确立为党的指导思想并\bk{写入党章}（胀气）；\\
  ②毛泽东在会上作题为《论联合政府》的书面报告和《愚公移山》闭幕词（七个小矮人联合起来愚公移山）；\\
  ③确定了党的政治路线：放手发动群众，壮大人民力量，在我党的领导下，打败日本侵略者，解放全国人民，建立一个新民主主义的中国；\\
  ④\kw{标志着我们党在政治上、思想上、组织上走向了成熟}（中共七大三成熟）：\\
  \hspace{1em}→\bk{在政治上}，党通过延安整风，使全党团结在毛泽东的旗帜下，实现了党的空前统一和团结；\\
  \hspace{1em}→\bk{在思想上}，党确立了毛泽东思想在全党的指导地位，把毛泽东思想写入了党章；\\
  \hspace{1em}→\bk{在组织上}，党形成了一支高举毛泽东旗帜的久经考验的政治家集团。\\
  \hspace{1em}党的七大在党的历史上具有极其重要的地位，为党后来不断从胜利走向胜利指明了正确方向、开辟了正确道路。
}

\zzsection{八、会议：1949 年以后}

\zzcard{会议顺口溜（续）}{%
  \zzkey{七届二中转城市，七届三中转经济，十一三中改革开放，十一六中通《决议》，十二大上中特提，十三阶段与路线，十四市场经济体，十五邓理论成指导，十六三个代表小康奋斗目标，十七科学发展观，十八全面小康新时代，十九新时代。}
}

\zzlongcard{历次会议要点}{%
  \zzsub{七届二中全会}
  \kw{“七届二中转城市，一方针两转变两务必”}\\
  \bk{两个务必}：务必继续保持谦虚、谨慎、不骄、不躁的作风，\kw{务必继续保持艰苦奋斗的作风}。\bk{人民政协政党立}。\\
  \zzsub{七届三中全会}
  \kw{“七届三中转经济”}：制定了党在建国初期的战略策略方针——\kw{不要四面出击}（“不要四”）。\\
  \zzsub{八大}
  矛盾定方针；在经济建设上，大会坚持既反保守又反冒进，即在综合平衡中\kw{稳步前进}的方针（四平八稳）。看到八大，想到\bk{经济建设和发展生产力}。\\
  \zzsub{十一届三中全会}
  \kw{“十一三中改革开放”}：会议\bk{重新确立}了马克思主义的\kw{思想路线、政治路线和组织路线}，恢复民主集中制，审查历史遗留问题和重要领导人功过是非问题。（重新确立三个路线，所以是“三中”）\\
  \zzsub{十一届六中全会}
  \kw{“十一六中通《决议》，社初提”}：①实质上初步提出了\bk{在中国建设什么样的社会主义和怎样建设社会主义}的问题；②《关于建国以来党的若干历史问题的决议》\kw{标志着党在指导思想上拨乱反正的胜利完成}。\\
  \hspace{1em}邓小平对《决议》的起草提出\bk{三条指导原则}：第一，\kw{确立毛泽东同志的历史地位，坚持和发展毛泽东思想}，这是\bk{最核心}的一条；第二，对建国 30 年来历史上的大事，要进行实事求是的分析，包括一些负责同志的功过是非，要作出公正的评价；第三，这个总结\kw{宜粗不宜细}，总结过去是为了引导大家团结一致向前看。\\
  \zzsub{十二大}
  \kw{“十二大上中特提”}：1982 年邓小平在党的十二大开幕词中明确提出：\kw{走自己的道路，建设有中国特色的社会主义}。从此，“中国特色社会主义”成为我们党的全部理论和实践创新的\bk{主题}。这一重大命题的提出\kw{具有里程碑意义}，是中国共产党人对科学社会主义发展的开创性贡献。\\
  \zzsub{十二届三中全会与十三大}
  \kw{“十三阶段与路线”}：十二届三中全会指出我国社会主义经济是\kw{在公有制基础上的有计划的商品经济}（有计划）。十三大提出\bk{社会主义初级阶段的基本路线}。\\
  \zzsub{十四大}
  \kw{“十四市场经济体”}：确立\bk{社会主义市场经济体制}。\\
  \zzsub{十五大}
  \kw{“十五邓理论成指导”}：首次明确提出了\kw{社会主义初级阶段的基本纲领}。\\
  \zzsub{十六大}
  \kw{“十六三个代表、小康奋斗目标”}。\\
  \zzsub{十七大}
  \kw{“十七科学发展观”}。\\
  \zzsub{十八大、十八届三中全会}
  \kw{“十八全面小康新时代”}。十八届三中全会：\kw{经济体制改革是全面深化改革的重点}，要使\bk{市场}在资源配置中起决定性作用和更好发挥\bk{政府}作用。\\
  \zzsub{十九大}
  \kw{“十九新时代”}。\\
  \zzsub{补充：几届几中全会一般做什么？}
  \bk{一中全会}——人事；\bk{二中全会}——人事；\bk{三中全会}——经济；\bk{四中全会}——重大决策部署；\bk{五中全会}——国民经济规划；\bk{六中全会}——意识形态、党的建设；\bk{七中全会}——承前启后。
}

\zzsection{九、建国后（1949—1956）}

\zzsub{建国初期三大外交政策}

\zzcard{建国初期三大外交政策}{%
  \kw{“打扫屋子倒炉灶”}\\
  \bk{一边倒}、\bk{另起炉灶}、\bk{打扫干净屋子再请客}。
}

\zzsub{个体经济}

\zzcard{个体经济}{%
  \kw{“十字路口有路标，像‘个’”}：新民主主义社会的五种经济成分中，处于\kw{十字路口}的经济成分是——\bk{个体经济}。
}

\zzsub{中国社会的阶级构成}

\zzcard{中国社会的阶级构成}{%
  \kw{“工农小民”}（无产阶级也称为工人阶级）\\
  \bk{工人阶级}（\kw{中国革命最基本的动力}）、\bk{农民阶级}、\bk{小资产阶级}、\bk{民族资产阶级}。\\
  \zzwarn{\bk{农民是中国革命的主力军}；\bk{农民问题是中国革命的基本问题}。}
}

\zzsub{中国进入社会主义社会最主要的标志}

\zzcard{中国进入社会主义社会最主要的标志}{%
  1956 年底，随着社会主义改造的完成，以\bk{生产资料公有制、按劳分配和计划经济体制}为特征的\kw{社会主义经济制度}建立起来，这是中国进入社会主义社会\kw{最主要的标志}。\\
  \zzwarn{经济基础决定上层建筑。}
}

\zzsub{多党合作和政治协商格局的形成}

\zzcard{多党合作和政治协商格局的形成}{%
  ①\bk{喊话——五一口号}：得到了热烈回应，揭开了序幕，奠定了基础；\\
  ②\bk{回应——《对时局的意见》}：各党派和无党派公开自愿接受共产党的领导；\\
  ③\bk{谈话——毛泽东谈话（1949 春）}：“积极参与，共同建设新中国”，\kw{标志着民主党派地位的根本变化}。
}

\zzsub{三大改造}

\zzcard{资本主义工商业改造三大步骤}{%
  \begin{zztable2}
    步骤 & 特点与性质 \\
    \midrule
    \bk{第一步：初级形式国家资本主义}（计划订货、统购包销、代工厂） & 国家与私人企业在\bk{企业外部}的合作；性质：\kw{社会主义萌芽性质} \\
    \bk{第二步：个别企业公私合营} & 生产资料归国家和资本家共同占有，\bk{公方占领导地位}；性质：\kw{半社会主义性质} \\
    \bk{第三步：全行业公私合营} & 企业变成\bk{社会主义国营企业}；性质：\kw{完全的社会主义性质} \\
  \end{zztable2}
  \zzwarn{第一、二步都是\bk{四马分肥}。}
}

\zzcard{三大改造的性质对照}{%
  \begin{zztable3}
    对象 & 萌芽 & 社会主义 \\
    \midrule
    农业 & \bk{互助组} & \bk{高级农业生产合作社} \\
    手工业 & \bk{手工业生产合作小组} & \bk{手工业生产合作社} \\
    资本主义工商业 & \bk{初级形式的国家资本主义} & \bk{全行业公私合营} \\
  \end{zztable3}
  中间形态：农业为\bk{初级农业生产合作社}（半社会主义）；手工业为\bk{手工业供销合作小组}；资本主义工商业为\bk{个别企业公私合营}。
}

\zzcard{三大改造的原则、方法与方针}{%
  \kw{“农原立典”}\\
  \bk{农业社会主义改造遵循的原则}：\kw{自愿互利、典型示范和国家帮助}。\\
  \bk{对个体手工业的社会主义改造的方法}：\kw{说服教育、典型示范和国家帮助}（说法）。\\
  \zzwarn{补充——对手工业的社会主义改造的\bk{方针}：\kw{积极领导、稳步前进}（针积极、针滴稳）。对\bk{农业}的社会主义改造也是这个方针。}
}

\zzsub{邓小平理论}

\zzcard{邓小平理论}{%
  ①深刻回答了\kw{什么是社会主义、怎样建设社会主义}这一根本问题；\\
  ②深刻揭示\bk{社会主义本质}；\\
  ③为我国改革开放和社会主义现代化建设提供了行动指南；\\
  ④是中国特色社会主义理论体系的\kw{奠基之作}。\\
  \zzwarn{辨析——\bk{邓小平理论}：什么是社会主义、怎样建设社会主义（看到邓小平，就想社会主义）；\bk{“三个代表”重要思想}：建设什么样的党、怎样建设党；\bk{科学发展观}：实现什么样的发展、怎样发展；\bk{习近平新时代中国特色社会主义思想}：新时代坚持和发展什么样的中国特色社会主义、怎样坚持和发展中国特色社会主义，建设什么样的社会主义现代化强国、怎样建设社会主义现代化强国，建设什么样的长期执政的马克思主义政党、怎样建设长期执政的马克思主义政党。}
}

\zzcard{南方谈话与宣言书}{%
  \bk{南方谈话}：从理论上深刻回答了长期困扰和束缚人们思想的许多重大问题，是\kw{把改革开放和现代化建设推向新阶段的又一个解放思想、实事求是的宣言书}，不仅对之后召开的党的十四大具有十分重要的指导作用，而且对中国整个社会主义现代化建设事业具有重大而深远的意义。\\[2pt]
  \zzsub{两个“宣言书”辨析}
  ①\bk{开辟中国特色社会主义道路、开创中国特色社会主义理论的宣言书}：\kw{《解放思想，实事求是，团结一致向前看》}；\\
  ②\bk{推进改革开放和社会主义现代化建设进入新阶段的又一个解放思想、实事求是的宣言书}：\kw{南方谈话}。
}

\zzsub{宪法}

\zzcard{两部宪法文件辨析}{%
  \begin{zztable2}
    《共同纲领》 & 《中华人民共和国宪法》 \\
    \midrule
    三大改造前，政协第一届全体会议通过 & 三大改造后，1954 年 9 月第一届全国人民代表大会第一次会议通过 \\
    成为中国人民的\bk{大宪章}，在一个时期内起着\bk{新中国临时宪法}的作用 & 是一部\bk{社会主义类型的宪法} \\
    \zzwarn{还\bk{不是}社会主义类型的宪法} & 体现了\bk{人民民主原则和社会主义原则} \\
  \end{zztable2}
}
